\kefalaio{ΜΚΔ και ΕΚΠ: μοιράζω και συναντιέμαι}{Κεφάλαιο Α.1.5, «Χαρακτήρες διαιρετότητας – ΜΚΔ – ΕΚΠ – Ανάλυση αριθμού σε γινόμενο πρώτων παραγόντων» (τρίτο μέρος)}

\begin{anagke}
Αυτή την εβδομάδα η κυρία Αλεξάνδρα βάζει δύο προβλήματα που μοιάζουν να έχουν αντίθετη ιδέα.

Πρώτο: το σχολείο φτιάχνει κήπο. Η Α’ τάξη έχει 150 φυτά, 90 γλάστρες και 60 σακούλες χώμα. Θέλουν να φτιάξουν ίδια κιτ για τα τμήματα, ώστε κάθε κιτ να έχει και από τα τρία είδη και να μη μείνει τίποτα. «Πόσα το πολύ ίδια κιτ γίνονται;» Ο Στέλιος λέει αμέσως «δέκα». Η Φανή ελέγχει: δέκα κιτ θέλουν 15 φυτά, 9 γλάστρες και 6 σακούλες το καθένα. Βγαίνει. «Μα δεν είπαμε πόσα το πολύ;» λέει η Ζωή.

Δεύτερο: ο θείος Λευτέρης και ο Ηλίας κάθονται το βράδυ στην ακτή και βλέπουν δύο φάρους. Ο ένας αναβοσβήνει κάθε 12 δευτερόλεπτα, ο άλλος κάθε 18. Μόλις άναψαν και οι δύο μαζί. «Σε πόσο ξανά θα ανάψουν μαζί;» ρωτάει ο θείος.

Στο πρώτο θέλουμε το μεγαλύτερο κοινό κομμάτι: μοιράζουμε. Στο δεύτερο θέλουμε την πρώτη φορά που συμπίπτουν: συναντιόμαστε. Και τα δύο λύνονται με τους ίδιους αριθμούς, διαιρέτες και πολλαπλάσια.
\end{anagke}

\begin{lexi}{ΜΚΔ και ΕΚΠ}{μέγιστος, ελάχιστος, κοινός}
Οι δύο συντομογραφίες κρύβουν τρεις λέξεις η καθεμία. \lexiword{Κοινός} είναι αυτός που ανήκει και στους δύο. \lexiword{Μέγιστος} είναι ο μεγαλύτερος που υπάρχει (το «μέγας» στον υπερθετικό), και \lexiword{ελάχιστος} ο μικρότερος που υπάρχει (το «ελαχύς»).

Το \lexiword{ΜΚΔ} είναι ο Μέγιστος Κοινός Διαιρέτης: ο μεγαλύτερος αριθμός που διαιρεί και τους δύο. Το \lexiword{ΕΚΠ} είναι το Ελάχιστο Κοινό Πολλαπλάσιο: το μικρότερο (μη μηδενικό) πολλαπλάσιο που είναι και των δύο.
\end{lexi}

\begin{metafrasi}
Ο μεγαλύτερος αριθμός που διαιρεί και το 12 και το 18. & $\MKD(12,\,18)=6$\\
Ο μικρότερος αριθμός (όχι 0) που είναι πολλαπλάσιο και του 12 και του 18. & $\EKP(12,\,18)=36$\\
Οι $\alpha$ και $\beta$ δεν έχουν κοινό διαιρέτη εκτός από το 1. & $\MKD(\alpha,\beta)=1$ (πρώτοι μεταξύ τους)\\
Ο $\alpha$ διαιρεί τον $\beta$. & $\MKD(\alpha,\beta)=\alpha$, $\EKP(\alpha,\beta)=\beta$\\
Το γινόμενο των δύο αριθμών. & $\MKD\cdot\EKP=\alpha\cdot\beta$\\
\end{metafrasi}

\tmima{Η θεωρία}

\begin{orismos}{ΜΚΔ}
Ο \textbf{Μέγιστος Κοινός Διαιρέτης} δύο ή περισσότερων φυσικών είναι ο μεγαλύτερος από τους κοινούς τους διαιρέτες. Γράφουμε $\MKD(\alpha,\beta)$.
\end{orismos}

\begin{orismos}{ΕΚΠ}
Το \textbf{Ελάχιστο Κοινό Πολλαπλάσιο} δύο ή περισσότερων φυσικών είναι το μικρότερο από τα κοινά τους πολλαπλάσια, εκτός από το 0 (που είναι πολλαπλάσιο όλων). Γράφουμε $\EKP(\alpha,\beta)$.
\end{orismos}

\tmimasmall{Δύο μέθοδοι}

\textbf{Με λίστες.} Γράφω τους διαιρέτες (για ΜΚΔ) ή τα πρώτα πολλαπλάσια (για ΕΚΠ) κάθε αριθμού και διαλέγω τον κοινό που ψάχνω. Είναι εύκολο για μικρούς αριθμούς.

\begin{kanonas}{ΜΚΔ και ΕΚΠ με ανάλυση}
\begin{enumerate}[leftmargin=1.6em,itemsep=0.2ex,topsep=2pt]
\item Αναλύω κάθε αριθμό σε γινόμενο πρώτων παραγόντων (με δυνάμεις).
\item \textbf{ΜΚΔ}: παίρνω μόνο τους \emph{κοινούς} πρώτους παράγοντες, καθένα με τον \emph{μικρότερο} εκθέτη που έχει.
\item \textbf{ΕΚΠ}: παίρνω \emph{όλους} τους πρώτους παράγοντες (κοινούς και μη), καθένα με τον \emph{μεγαλύτερο} εκθέτη που έχει.
\item Πολλαπλασιάζω τους παράγοντες που πήρα.
\end{enumerate}
\end{kanonas}

Η ιδέα: το ΜΚΔ πρέπει να χωράει και στους δύο αριθμούς, γι' αυτό παίρνει όσα έχουν και οι δύο, το λιγότερο. Το ΕΚΠ πρέπει να χωράει και τους δύο, γι' αυτό παίρνει όλα όσα χρειάζεται ο καθένας, το περισσότερο.

\tmimasmall{Πρώτοι μεταξύ τους και μια σχέση}

Δύο αριθμοί λέγονται \textbf{πρώτοι μεταξύ τους} αν ο μοναδικός κοινός διαιρέτης τους είναι το 1, δηλαδή αν $\MKD(\alpha,\beta)=1$. Δεν χρειάζεται να είναι οι ίδιοι πρώτοι: οι 8 και 15 είναι σύνθετοι αλλά πρώτοι μεταξύ τους. Τότε το ΕΚΠ τους είναι το γινόμενό τους: $\EKP(8,\,15)=120$.

Για δύο φυσικούς ισχύει πάντα
$$\MKD(\alpha,\beta)\cdot\EKP(\alpha,\beta)=\alpha\cdot\beta.$$
Για τα 12 και 18: $6\cdot 36=216=12\cdot 18$.

\begin{prosoxi}
\begin{itemize}[leftmargin=1.2em,itemsep=0.3ex,topsep=2pt]
\item Το ΜΚΔ είναι πάντα μικρότερο ή ίσο από τον μικρότερο αριθμό. Το ΕΚΠ είναι πάντα μεγαλύτερο ή ίσο από τον μεγαλύτερο. Αν βγάλεις ΜΚΔ μεγαλύτερο από έναν από τους αριθμούς, κάτι πήγε στραβά.
\item Στο ΜΚΔ προσοχή: παίρνω μόνο τους \emph{κοινούς} παράγοντες. Στο ΕΚΠ παίρνω \emph{όλους}.
\item Όταν ένας αριθμός διαιρεί τον άλλον, τα πράγματα είναι εύκολα: $\MKD(7,42)=7$ και $\EKP(7,42)=42$.
\end{itemize}
\end{prosoxi}

\tmima{Λυμένα παραδείγματα}

\paradeigma{Οι φάροι}
\tixero Ο ένας φάρος αναβοσβήνει κάθε 12 δευτερόλεπτα, ο άλλος κάθε 18. Άναψαν μαζί τώρα.\par
\tithelo Σε πόσα δευτερόλεπτα θα ανάψουν ξανά μαζί.\par
\possindeo Ο πρώτος ανάβει στα δευτερόλεπτα $12, 24, 36, 48, \dots$ και ο δεύτερος στα $18, 36, 54, \dots$ Η πρώτη φορά που συμπίπτουν είναι το 36. Δηλαδή ψάχνω το $\EKP(12,18)$.\par
Με ανάλυση: $12=2^2\cdot 3$ και $18=2\cdot 3^2$. Παίρνω όλους τους παράγοντες με τον μεγαλύτερο εκθέτη: $2^2\cdot 3^2=36$.\par
\apantisi Σε 36 δευτερόλεπτα.\par
\elegxos Το 36 είναι πολλαπλάσιο και του 12 ($12\cdot 3$) και του 18 ($18\cdot 2$), και κανένας μικρότερος αριθμός δεν είναι (δες τις λίστες).

\paradeigma{Τα κιτ του κήπου}
\tixero 150 φυτά, 90 γλάστρες και 60 σακούλες χώμα. Τα κιτ πρέπει να είναι ίδια και να μη μείνει τίποτα.\par
\tithelo Πόσα το πολύ κιτ γίνονται και τι έχει το καθένα.\par
\possindeo Ο αριθμός των κιτ πρέπει να διαιρεί και το 150 και το 90 και το 60. Ψάχνω το μεγαλύτερο κοινό διαιρέτη. $150=2\cdot 3\cdot 5^2$, $90=2\cdot 3^2\cdot 5$, $60=2^2\cdot 3\cdot 5$. Κοινοί παράγοντες: 2, 3 και 5, με τον μικρότερο εκθέτη τον 1 σε κάθε περίπτωση. Άρα $\MKD=2\cdot 3\cdot 5=30$.\par
Κάθε κιτ έχει: $150:30=5$ φυτά, $90:30=3$ γλάστρες, $60:30=2$ σακούλες.\par
\apantisi 30 κιτ, με 5 φυτά, 3 γλάστρες και 2 σακούλες το καθένα.\par
\elegxos Τα 10 κιτ του Στέλιου βγαίνουν, αλλά το 30 είναι κοινός διαιρέτης και μεγαλύτερος από το 10. Αν δοκιμάσουμε μεγαλύτερο αριθμό, π.χ. 60 κιτ, το 90 δεν διαιρείται με το 60 ($90:60$ δεν είναι ακέραιο).

\paradeigma{ΜΚΔ και ΕΚΠ τριών αριθμών}
\tixero Οι αριθμοί 2.520, 2.940 και 3.780.\par
\tithelo Το ΜΚΔ και το ΕΚΠ τους.\par
\possindeo Από το προηγούμενο κεφάλαιο: $2.520=2^3\cdot 3^2\cdot 5\cdot 7$. Με τον ίδιο τρόπο $2.940=2^2\cdot 3\cdot 5\cdot 7^2$ και $3.780=2^2\cdot 3^3\cdot 5\cdot 7$.\par
ΜΚΔ: οι παράγοντες 2, 3, 5, 7 υπάρχουν και στους τρεις. Παίρνω τον μικρότερο εκθέτη του καθενός: $2^2\cdot 3\cdot 5\cdot 7=420$.\par
ΕΚΠ: παίρνω τον μεγαλύτερο εκθέτη του καθενός: $2^3\cdot 3^3\cdot 5\cdot 7^2=8\cdot 27\cdot 5\cdot 49=52.920$.\par
\apantisi $\MKD=420$ και $\EKP=52.920$.\par
\elegxos Το 420 διαιρεί τον καθένα: $2.520:420=6$, $2.940:420=7$, $3.780:420=9$. Και το 52.920 είναι πολλαπλάσιο του καθενός: $52.920:2.520=21$.

\tmima{Ασκήσεις}

\zoni{Το κατάλαβα}{εφαρμογή του κανόνα}

\ex Βρες το ΜΚΔ γράφοντας τους διαιρέτες: α) 12 και 30 \quad β) 18 και 24 \quad γ) 20 και 45 \quad δ) 16 και 27.

\ex Βρες το ΕΚΠ γράφοντας τα πρώτα πολλαπλάσια: α) 4 και 6 \quad β) 6 και 15 \quad γ) 8 και 12 \quad δ) 5 και 7.

\ex Βρες το ΜΚΔ και το ΕΚΠ με ανάλυση σε πρώτους παράγοντες: α) 24 και 36 \quad β) 60 και 84 \quad γ) 90 και 150 \quad δ) 48 και 180.

\ex Βρες το ΜΚΔ και το ΕΚΠ των τριών αριθμών 12, 18 και 30.

\ex Ποια από τα ζευγάρια είναι πρώτοι μεταξύ τους; α) 8 και 15 \quad β) 12 και 35 \quad γ) 21 και 35 \quad δ) 9 και 28 \quad ε) 14 και 49 \quad στ) 1 και 99.

\ex Βρες το ΜΚΔ και το ΕΚΠ των: α) 7 και 42 \quad β) 15 και 60 \quad γ) 25 και 100. Τι παρατηρείς όταν ο ένας αριθμός διαιρεί τον άλλον;

\ex Σωστό ή λάθος; α) $\MKD(12,18)=36$. \quad β) $\EKP(4,6)=24$. \quad γ) Το ΜΚΔ δύο διαφορετικών πρώτων είναι 1. \quad δ) Το ΕΚΠ δύο αριθμών δεν είναι ποτέ μικρότερο από τον μεγαλύτερο. \quad ε) Το ΜΚΔ δεν είναι ποτέ μεγαλύτερο από τον μικρότερο. \quad στ) Αν δύο αριθμοί είναι πρώτοι μεταξύ τους, το ΕΚΠ τους είναι το γινόμενό τους.

\zoni{Το χρησιμοποιώ}{προβλήματα από τη ζωή}

\ex \textbf{Οι φάροι.} Δύο άλλοι φάροι αναβοσβήνουν ο ένας κάθε 15 δευτερόλεπτα και ο άλλος κάθε 20. Άναψαν μαζί τώρα. Πότε θα ξανανάψουν μαζί; Πόσες φορές θα συμπέσουν στα επόμενα 3 λεπτά, μετρώντας και το τέλος του τρίτου λεπτού;

\ex \textbf{Τα καλάθια της γιαγιάς.} Η γιαγιά Ρήνη έχει 84 μήλα και 126 πορτοκάλια και θέλει να φτιάξει ίδια καλάθια, με και από τα δύο φρούτα και χωρίς να περισσέψει τίποτα. Πόσα το πολύ καλάθια μπορεί να φτιάξει; Τι έχει το καθένα;

\ex \textbf{Τα πλακάκια του θείου.} Η βεράντα του θείου Λευτέρη είναι ορθογώνια, 180 εκ. επί 150 εκ. Θέλει να την καλύψει με τετράγωνα πλακάκια ίδιου μεγέθους, χωρίς να κόψει κανένα. Ποια είναι η μεγαλύτερη δυνατή πλευρά του πλακιδίου; Πόσα πλακάκια θα χρειαστεί;

\ex \textbf{Τα λεωφορεία.} Δύο γραμμές λεωφορείων ξεκινούν μαζί από τη στάση στις 7:00. Η μία περνάει κάθε 15 λεπτά και η άλλη κάθε 20. Πότε θα ξανακαταλήξουν μαζί; Πόσες φορές θα βρεθούν μαζί από τις 7:00 ως τις 10:00 (μετρώντας και τις δύο ώρες);

\ex \textbf{Το πότισμα.} Η Μυρτώ ποτίζει τις γλάστρες κάθε 4 μέρες, ο Ηλίας κάθε 6 και η Ζωή κάθε 9. Σήμερα τις πότισαν και οι τρεις. Σε πόσες μέρες θα τις ποτίσουν ξανά και οι τρεις την ίδια μέρα;

\zoni{Το σκέφτομαι}{μια ερώτηση που πάει πιο μακριά}

\ex Για τα ζευγάρια 12 και 18, 8 και 20, 9 και 15, 14 και 21 βρες το ΜΚΔ, το ΕΚΠ και το γινόμενο των δύο αριθμών. Σύγκρινε το γινόμενο ΜΚΔ·ΕΚΠ με το γινόμενο των αριθμών. Τι ανακαλύπτεις;

\ex Η Ζωή παρατηρεί ότι $48=18\cdot 2+12$, και ότι $\MKD(48,18)=6=\MKD(18,12)$. Δοκίμασε το ίδιο με τα 84 και 30: γράψε την Ευκλείδεια διαίρεση, και σύγκρινε το $\MKD(84,30)$ με το $\MKD(30,\upsilon)$. Αν συνεχίσεις μέχρι να βρεις υπόλοιπο 0, τι δείχνει ο τελευταίος διαιρέτης; Δοκίμασε με τα 252 και 105. (Αυτή είναι η μέθοδος του Ευκλείδη.)

\ex Γιατί, όταν δύο αριθμοί είναι πρώτοι μεταξύ τους, το ΕΚΠ είναι το γινόμενό τους; Εξήγησε με τα 8 και 15, χρησιμοποιώντας τις αναλύσεις τους.

\ex Βρες δύο αριθμούς που να έχουν ΜΚΔ 6 και ΕΚΠ 36. Υπάρχει και δεύτερο ζευγάρι; (Βοήθεια: το γινόμενο των δύο πρέπει να είναι $6\cdot 36$.)

\ex Μπορεί το ΜΚΔ δύο \emph{διαφορετικών} αριθμών να είναι ίσο με το ΕΚΠ τους; Εξήγησε.
